\documentclass[11pt,a4paper]{article}
\usepackage[margin=1in]{geometry}
\usepackage{amsmath,amssymb}
\usepackage{booktabs}
\usepackage{enumitem}
\usepackage{titlesec}
\usepackage{parskip}
\usepackage{multirow}
\usepackage{array}
\usepackage{fancyhdr}
\usepackage[hidelinks]{hyperref}
\usepackage{graphicx}
\usepackage{siunitx}
\sisetup{detect-weight=true, detect-family=true, detect-shape=true}
\usepackage{caption}
\usepackage{subcaption}

\usepackage{biblatex}
\addbibresource{sample.bib}

\pagestyle{fancy}
\fancyhf{}
\rhead{\small Advanced Machine Learning --- Final Report}
\lhead{\small Innopolis University, 2026}
\cfoot{\thepage}

\titleformat{\section}{\large\bfseries}{}{0em}{}[\titlerule]
\titleformat{\subsection}{\normalsize\bfseries}{}{0em}{}

\setlength{\parskip}{6pt}

\begin{document}

\begin{center}
    {\LARGE \textbf{Analyzing Inductive Bias in Biomedical Retrieval:\\[4pt]
    A Multi-Model, Multi-Dataset Study}}\\[10pt]
    {\large Final Report}\\[6pt]
    {\normalsize Dmitrii Naumov \quad|\quad Innopolis University \quad|\quad Advanced ML, Spring 2026}\\[4pt]
    {\small \today}
\end{center}

\vspace{4pt}
\hrule
\vspace{10pt}

\begin{abstract}
\noindent
This report studies how the inductive bias of a retriever shapes
performance on biomedical retrieval. Six retrievers (BM25, MiniLM,
BGE-small, E5-small, SPLADE, MedCPT) are compared on five BEIR-style
benchmarks. BGE-small is the strongest single system on NFCorpus
(nDCG@10 = 0.3388); E5-small wins at nDCG@1 (0.4231). MedCPT
collapses to 0.1332 on ArguAna, a 0.228 loss against BM25. A
general-domain BGE cross-encoder reranker hurts on four of five
datasets; the in-domain MedCPT cross-encoder improves NFCorpus by
0.027. For per-query routing on NFCorpus, the oracle reaches 0.4097,
$+0.066$ above the static BM25+BGE-small hybrid (0.3434); the learned
routers tested here do not close that gap, so the routing signal
exists but is hard to predict from query features alone. Two rows are
reported as caveats rather than claims: BioASQ-subset is a
degenerate-corpus stress case (zero distractors on the run that
produced these numbers), and the PubMedQA-labeled section is a
saturation sanity check on a tiny candidate pool. The findings support
No Free Lunch: no single retriever or reranker dominates, and
biomedical gains come from domain-aligned reranking rather than from
stronger general-domain models.
\end{abstract}

\vspace{4pt}
\hrule
\vspace{10pt}

% -------------------------------------------------------
\section{Problem and Research Question}

Information retrieval is the first stage of a retrieval-augmented
generation pipeline. The choice of retriever decides which documents
reach the generator, and therefore how well the answer is grounded in
evidence. Every retrieval method encodes implicit assumptions about
relevance.

\textbf{BM25}~\cite{Robertson2009BM25} treats relevance as exact lexical
overlap, weighted by term frequency and inverse document frequency.
\textbf{General-domain dense retrieval} (\texttt{all-MiniLM-L6-v2}~\cite{MiniLM},
BGE~\cite{Xiao2023BGE}, E5~\cite{Wang2022E5}) measures relevance as cosine
similarity in a learned embedding space. \textbf{Learned sparse retrieval}
(SPLADE~\cite{Formal2021SPLADE}) combines lexical precision with
contrastively-trained term expansion. \textbf{Domain-adapted dense
retrieval} (MedCPT~\cite{Jin2023MedCPT}) keeps the bi-encoder architecture
but trains it on biomedical click-through data.

\paragraph{Scope.}
The main analysis evaluates retrieval quality. A single generation
section (Section~\ref{sec:downstream}) is added so that the RAG framing
of the report is grounded in a downstream task with gold answers, but
generation is not the focus.

\paragraph{Research question.}
\emph{Does the inductive bias of a retrieval method produce a stable and
explainable performance difference on biomedical retrieval, or does the
observed advantage depend on the dataset, query sample, or evaluation
protocol?} This question connects to three course concepts:
(1) \textbf{No Free Lunch} (no inductive bias is universally optimal);
(2) \textbf{empirical vs.\ true risk} (how well the measured nDCG@10
on the test set approximates the expected nDCG@10 on the full query
distribution); and (3) \textbf{stability} (whether the observed
ranking is consistent across query subsets and datasets).

% -------------------------------------------------------
\section{Datasets}
\label{sec:datasets}

The study uses five BEIR-style benchmarks from biomedical, scientific
and argumentation domains. Their statistics are summarised in
Table~\ref{tab:datasets}.

\begin{table}[ht]
\centering
\caption{Datasets used in this study. NFCorpus, TREC-COVID, SciFact and
ArguAna use the BEIR-distributed splits. BioASQ is evaluated on a
deterministic subset described below.}
\label{tab:datasets}
\small
\begin{tabular}{lrrll}
\toprule
\textbf{Dataset} & \textbf{\# Docs} & \textbf{\# Test Q} & \textbf{Domain} & \textbf{Query style} \\
\midrule
NFCorpus      & 3{,}633   & 323     & Biomedical          & Plain-language health questions \\
BioASQ-subset & 28{,}001  & 500     & Biomedical (PubMed) & Expert factoid / yes-no questions \\
TREC-COVID    & 171{,}332 & 50      & Biomedical (COVID)  & Short topic statements \\
SciFact       & 5{,}183   & 300     & Scientific claims   & Atomic scientific claims \\
ArguAna       & 8{,}674   & 1{,}406 & Argumentation       & Full-sentence counter-arguments \\
\bottomrule
\end{tabular}
\end{table}

\paragraph{NFCorpus version.}
The original NFCorpus release~\cite{boteva2016} contains 9{,}964
documents. The BEIR-distributed version used here keeps only 3{,}633
documents (those that appear in at least one relevance judgement). All
NFCorpus numbers in this report refer to the BEIR version and cannot be
compared with results reported on the full 9{,}964-document
corpus.

\paragraph{Computational environment.}
The project was developed locally on a MacBook M4 Pro, but several
models (in particular SPLADE and MedCPT) were too slow to encode on
CPU. The final experiments were therefore run on Google Colab with a
Tesla T4 GPU. BM25 retrieval and FAISS search use the CPU. Latency
numbers in Section~\ref{subsec:efficiency} reflect this environment.

\paragraph{BioASQ subset.}
The full BioASQ-BEIR corpus~\cite{Tsatsaronis2015BioASQ} is too large
for a six-retriever sweep on a Colab T4 in a single pass. The notebook
tries three loaders in sequence (a local BEIR-format directory, the
\texttt{BeIR/bioasq} mirror, and a public mini-BioASQ mirror); the first
one that loads correctly is used. The deterministic subset is then
built with seed 42 from (i) all documents that appear in at least one
test qrel and (ii) up to 10{,}000 sampled distractors. The actual run
that produced the numbers below loaded the mini-BioASQ mirror, whose
corpus pool is exhausted by the qrel union itself: the resulting subset
contains \emph{28{,}001 documents (qrel-union, zero distractors)}
rather than the planned 50{,}000. The 500-query sample on top of the
subset was kept as planned. The implications of the no-distractor
setup are discussed when the BioASQ row is interpreted in
Section~\ref{subsec:multi}.

\paragraph{Why biomedical datasets are the focus.}
NFCorpus, BioASQ-subset and TREC-COVID form the biomedical family.
SciFact is a scientific but not strictly biomedical control, and
ArguAna is a non-biomedical control. The three biomedical sets test
whether the general-domain pre-training of MiniLM, BGE and E5 misses
technical terminology that a domain-adapted encoder like MedCPT can
capture. SciFact and ArguAna check how much of any biomedical
alignment transfers out of domain.

\subsection{Primary metric}

The primary metric is nDCG@10:
\begin{equation}
    \text{nDCG@}k = \frac{\text{DCG@}k}{\text{IDCG@}k},
    \qquad
    \text{DCG@}k = \sum_{i=1}^{k} \frac{2^{r_i} - 1}{\log_2(i+1)}.
\end{equation}
Per-query loss is $\ell(h,q) = 1 - \text{nDCG@10}(h,q) \in [0,1]$.
Empirical risk is $\hat{R}(h) = \frac{1}{n}\sum_{q}\ell(h,q)$.
MAP@10, Recall@10 and P@10 are reported as secondary metrics.

% -------------------------------------------------------
\section{Retrieval Methods}

Six retrievers are evaluated, covering four inductive-bias classes
(Table~\ref{tab:retrievers}). All systems retrieve the top-100
candidates and metrics are computed at the top-10 cut-off.

\begin{table}[ht]
\centering
\caption{Retrievers compared in this study.}
\label{tab:retrievers}
\small
\begin{tabular}{llll}
\toprule
\textbf{System} & \textbf{Type} & \textbf{Bias class} & \textbf{Model / parameters} \\
\midrule
BM25              & Lexical sparse  & Exact term overlap                & \texttt{bm25s}, $k_1\!=\!1.5, b\!=\!0.75$ \\
Dense (MiniLM)    & Bi-encoder      & General-domain semantic similarity & \texttt{all-MiniLM-L6-v2} (22\,M) \\
BGE-small         & Bi-encoder      & Contrastively-trained semantic     & \texttt{BAAI/bge-small-en-v1.5} (33\,M) \\
E5-small          & Bi-encoder      & Weakly-supervised semantic         & \texttt{intfloat/e5-small-v2} (33\,M) \\
SPLADE            & Learned sparse  & Supervised term expansion          & \texttt{splade-cocondenser-ensembledistil} \\
MedCPT            & Bi-encoder      & Domain-adapted biomedical          & \texttt{ncbi/MedCPT-\{Query,Article\}-Encoder} \\
\bottomrule
\end{tabular}
\end{table}

BM25 follows the Lucene IDF variant. Its inductive bias is exact
lexical overlap; it cannot bridge synonyms or paraphrases. The three
dense bi-encoders (MiniLM, BGE-small, E5-small) encode query and
document into $\mathbb{R}^{384}$ and use cosine similarity (computed as
the inner product of L2-normalised embeddings). E5 uses the documented
\texttt{query:} and \texttt{passage:} prefixes. SPLADE applies a
log--ReLU activation on BERT MLM logits and pools across tokens, which
gives a sparse vocabulary-sized vector per text; retrieval is sparse
matrix multiplication. MedCPT is a dual 768-dimensional CLS-token
bi-encoder trained on 255 million PubMed click-through pairs; its
architecture is similar to MiniLM/BGE/E5, but its pre-training
distribution matches biomedical retrieval much more closely.

\paragraph{Hybrid.}
A linear score fusion of BM25 and a dense system after per-method
min--max normalisation is also evaluated:
\begin{equation}
    \text{score}_{\text{hybrid}}(q,d) =
    \alpha \cdot \hat{s}_{\text{BM25}}(q,d) + (1-\alpha) \cdot \hat{s}_{\text{Dense}}(q,d),
    \qquad \alpha \in [0,1].
    \label{eq:hybrid}
\end{equation}
The weight $\alpha$ is tuned on the NFCorpus dev split. Reciprocal Rank
Fusion (RRF, $k=60$) is also reported as a parameter-free baseline.

% -------------------------------------------------------
\section{Results}

Phase-A first-stage results were produced on Google Colab with a Tesla
T4 GPU for dense encoder inference and CPU for BM25 and FAISS retrieval.
The Phase-B additions (BioASQ-subset retrieval, BGE cross-encoder
reranking, the per-query router and the downstream MIRAGE section) were
produced by a single consolidated Colab driver that runs end-to-end on
one T4 session. Earlier next-stage analyses (paired-bootstrap CIs,
the $\Delta(q)$ regression, and the stratification tables in
Section~\ref{sec:theory}) are reused from a previous run.

\subsection{Retrieval quality on NFCorpus}
\label{subsec:nf_quality}

Table~\ref{tab:nfcorpus_6way} reports the six-way comparison on the
323-query NFCorpus test split. The ranking depends on the cut-off:
E5-small is best at nDCG@1, BGE-small is best at nDCG@3, nDCG@5 and
nDCG@10. MiniLM-Dense is the weakest of the four bi-encoders at
nDCG@1 but moves above BM25 from rank~3 onwards.

\begin{table}[ht]
    \centering
    \caption{Retrieval quality on the NFCorpus test split (323 queries).
    The best result per metric is in bold. SPLADE and MedCPT are reported
    only at the top-10 cut-off because the original logging for those
    two systems was limited to nDCG@10, MAP@10, Recall@10 and P@10.}
    \label{tab:nfcorpus_6way}
    \small
    \begin{tabular}{lSSSSSS}
\toprule
\textbf{Metric} & \textbf{BM25} & \textbf{Dense} & \textbf{BGE-small} & \textbf{E5-small} & \textbf{SPLADE} & \textbf{MedCPT} \\
\midrule
nDCG@1 & 0.4087 & 0.3911 & 0.4138 & \textbf{0.4231} & {---} & {---} \\
nDCG@3 & 0.3500 & 0.3605 & \textbf{0.3920} & 0.3673 & {---} & {---} \\
nDCG@5 & 0.3264 & 0.3400 & \textbf{0.3715} & 0.3513 & {---} & {---} \\
nDCG@10 & 0.2940 & 0.3173 & \textbf{0.3388} & 0.3267 & 0.3327 & 0.3246 \\
\midrule
MAP@10 & 0.2081 & 0.2277 & \textbf{0.2456} & 0.2333 & 0.2439 & 0.2324 \\
Recall@10 & 0.1373 & 0.1550 & 0.1583 & 0.1595 & 0.1593 & \textbf{0.1599} \\
P@10 & 0.2074 & 0.2433 & \textbf{0.2508} & 0.2390 & 0.2390 & 0.2424 \\
\midrule
Empirical risk $\hat{R}$ & 0.7060 & 0.6827 & \textbf{0.6612} & 0.6733 & 0.6673 & 0.6754 \\
Loss variance & 0.0967 & 0.0970 & 0.1018 & 0.1006 & 0.1014 & 0.1000 \\
\bottomrule
\end{tabular}

\end{table}

The nDCG@10 ranking is BGE-small (0.3388) $>$ SPLADE (0.3327) $>$
E5-small (0.3267) $>$ MedCPT (0.3246) $>$ MiniLM-Dense (0.3173) $>$ BM25
(0.2940). MedCPT is below BGE-small even though it is biomedical.
A likely reason is that NFCorpus questions are written in
plain-language consumer-health style, while MedCPT was trained on
PubMed search clicks, where the vocabulary and intent are different.
At nDCG@1 the picture changes: E5-small (0.4231), BGE-small (0.4138)
and BM25 (0.4087) are within 0.015 of each other; the bi-encoder
advantage grows with the cut-off because dense models surface
semantically related but lexically distant documents in the lower
ranks.

\subsection{Multi-dataset generalisation}
\label{subsec:multi}

The biomedical claim of the report rests on three biomedical
benchmarks: NFCorpus, TREC-COVID and BioASQ-subset. SciFact (scientific
claims) and ArguAna (argumentation) are kept as out-of-domain controls,
not as substitutes for BioASQ. SPLADE and MedCPT are skipped on
TREC-COVID because encoding the 171{,}332-document corpus exceeded the
available GPU budget on Colab. The NFCorpus row in
Table~\ref{tab:multi_dataset} uses the same canonical numbers as
Table~\ref{tab:nfcorpus_6way}.

\begin{table}[ht]
    \centering
    \caption{nDCG@10 across all retrievers and datasets. The best result
    per row is in bold. Cells marked ``--'' on TREC-COVID denote runs
    skipped due to the available compute budget. The dagger ($\dagger$)
    on BioASQ-subset marks a degenerate-corpus row that is kept as a
    stress case and not used as biomedical evidence
    (see Section~\ref{subsec:multi}).}
    \label{tab:multi_dataset}
    \small
    \begin{tabular}{lSSSSSS}
\toprule
\textbf{Dataset} & \textbf{BM25} & \textbf{Dense} & \textbf{BGE-small} & \textbf{E5-small} & \textbf{SPLADE} & \textbf{MedCPT} \\
\midrule
NFCorpus & 0.2940 & 0.3173 & \textbf{0.3388} & 0.3267 & 0.3327 & 0.3246 \\
BioASQ-subset$^{\dagger}$ & \textbf{0.7544} & 0.5392 & 0.6836 & 0.6658 & 0.6733 & 0.6101 \\
TREC-COVID & 0.5759 & 0.4539 & 0.6448 & \textbf{0.7238} & {---} & {---} \\
SciFact & 0.6617 & 0.6451 & \textbf{0.7200} & 0.6875 & 0.6327 & 0.7103 \\
ArguAna & 0.3606 & 0.3698 & \textbf{0.4287} & 0.3100 & 0.3942 & 0.1332 \\
\bottomrule
\end{tabular}

% $^{\dagger}$ BioASQ-subset is a stress case, not biomedical evidence. See §\ref{subsec:multi}.

\end{table}

\paragraph{BioASQ-subset caveat.}
The BioASQ row must be interpreted with care. The mini-BioASQ source
that succeeded on this run produced a 28{,}001-document corpus that is
identical to the qrel union, with zero distractors. Every document in
the index is therefore relevant to at least one query. The retrieval
task becomes artificially easy and biases the comparison towards
exact-match methods: BM25 nDCG@10 = 0.7544 here, against a published
full-corpus baseline of about 0.46 ($+0.30$ inflation). Every learned
retriever loses to BM25 by margins that are at least partly driven by
the same effect.

\textbf{The BioASQ-subset row is therefore not used as evidence about
biomedical retrieval anywhere in this report.} It is kept in
Table~\ref{tab:multi_dataset} (and the corresponding rows in
Table~\ref{tab:reranker}) as a stress case, so that the reader can see
how the same retrievers and rerankers behave when the corpus is
degenerate. All biomedical claims in this report rest on NFCorpus and
TREC-COVID; SciFact and ArguAna are kept as out-of-domain controls.
Re-running BioASQ with a real distractor corpus (e.g.\ the gated
official BioASQ-BEIR release, or PubMed abstracts as distractors) is
listed in Section~\ref{sec:limitations} as the natural fix and is left
to future work.

\paragraph{Mean differences.}
Mean nDCG@10 differences against BM25 across the five datasets are:
BGE-small $+0.034$, E5-small $+0.013$, SPLADE $-0.009$, MiniLM-Dense
$-0.064$, MedCPT $-0.073$. The BioASQ row pushes the means of every
learned retriever down because of the no-distractor issue. If BioASQ
is excluded, the means become $+0.060$, $+0.039$, $+0.013$, $-0.027$
and $-0.052$ respectively. The ordering is the same, but BGE-small
becomes positive against BM25 again on the four non-degenerate
datasets.

Three patterns from the four non-degenerate datasets support the No
Free Lunch reading. MiniLM-Dense beats BM25 on NFCorpus and ArguAna
but loses on TREC-COVID by $-0.119$, where rare COVID-era biomedical
vocabulary needs exact matching. MedCPT outperforms BM25 on NFCorpus
and SciFact but drops to 0.1332 on ArguAna, a loss of $-0.227$. This
is the largest undisputed negative gap in the study; biomedical
pre-training actively hurts on a non-biomedical argumentation task.
BGE-small is the most stable single system: it wins on NFCorpus,
SciFact and ArguAna, and is second on TREC-COVID. Its average gap of
$+0.060$ against BM25 over the four non-degenerate datasets is the
largest of any learned system.

A paired bootstrap on the 500-query BioASQ-subset sample (10{,}000
resamples) gives MedCPT $-$ BM25 $= -0.1443$, 95\% CI $[-0.1724,
-0.1167]$, $P(\text{MedCPT}>\text{BM25})=0$. The CI is narrow, but
the absolute magnitude is inflated by the no-distractor index.

\subsection{Cross-encoder reranking}
\label{subsec:reranker}

A cross-encoder reranking track is added on top of the strongest
first-stage retriever for each dataset. The general-domain reranker is
\texttt{BAAI/bge-reranker-base}~\cite{Xiao2023BGE}, a BERT-style model
that scores $(q, d)$ pairs jointly. An in-domain biomedical reranker,
\texttt{ncbi/MedCPT-Cross-Encoder}~\cite{Jin2023MedCPT}, is also
evaluated on NFCorpus and BioASQ-subset. Both rerankers see the
top-100 candidates of the first stage. Results are reported in
Table~\ref{tab:reranker}.

\begin{table}[ht]
    \centering
    \caption{Cross-encoder (BGE-reranker-base) and biomedical
    cross-encoder (MedCPT-CE) reranking over first-stage top-100
    candidates. The first-stage column reuses the strongest single
    retriever from Section~\ref{subsec:multi}; the Delta column is the
    change in mean nDCG@10 with respect to the first stage.}
    \label{tab:reranker}
    \small
    \begin{tabular}{llSSS}
\toprule
\textbf{Dataset} & \textbf{First-stage} & \textbf{First-stage nDCG@10} & \textbf{Reranked nDCG@10} & \textbf{$\Delta$} \\
\midrule
NFCorpus & BGE-small + BGE-reranker@100 & 0.3393 & 0.3007 & -0.0386 \\
NFCorpus & BM25 + BGE-reranker@100 & 0.3076 & \textbf{0.3082} & +0.0006 \\
BioASQ-subset$^{\dagger}$ & BM25 + BGE-reranker@100 & 0.7544 & 0.7080 & -0.0464 \\
TREC-COVID & E5-small + BGE-reranker@100 & 0.7238 & \textbf{0.7540} & +0.0302 \\
SciFact & BGE-small + BGE-reranker@100 & 0.7200 & 0.7091 & -0.0109 \\
ArguAna & BGE-small + BGE-reranker@100 & 0.4287 & 0.1860 & -0.2428 \\
\midrule
NFCorpus & BGE-small + MedCPT-CE@100 & 0.3393 & \textbf{0.3664} & +0.0271 \\
BioASQ-subset$^{\dagger}$ & BM25 + MedCPT-CE@100 & 0.7544 & \textbf{0.8057} & +0.0513 \\
\bottomrule
\end{tabular}

% $^{\dagger}$ degenerate corpus, see BioASQ caveat.

\end{table}

\paragraph{Calibration note.}
First-stage NFCorpus values in Tables~\ref{tab:reranker}
and~\ref{tab:router} are remeasured inside the consolidated Phase-B
pipeline and differ slightly from the canonical Phase-A values
(BM25 0.2940 vs.\ 0.3076, BGE-small 0.3388 vs.\ 0.3393). The drift is
caused by a different BM25 tokenizer and is well below the
paired-bootstrap CI half-width of 0.034
(Section~\ref{subsec:paired_bootstrap}). The $\Delta$ column is
calibration-invariant.

\paragraph{Reading the results.}
The cross-encoder track gives three observations. BGE-reranker-base
does not help on most datasets. Of the five first-stage / dataset
pairs only TREC-COVID gains ($+0.0302$); NFCorpus BGE-small drops by
$-0.0386$, SciFact by $-0.0109$, ArguAna by $-0.2428$, and BioASQ by
$-0.0464$ on the degenerate-corpus index (read as a stress case).
A stronger general-domain reranker is not automatically a stronger
pipeline.

The in-domain MedCPT cross-encoder gives a clean positive delta on
NFCorpus ($+0.0271$). The BioASQ row also moves up by $+0.0513$, but
this is on the same degenerate corpus as the first-stage BioASQ row
and is reported only for completeness; it is not a clean biomedical
reranker claim. The clean biomedical reranker result in this report is
therefore the NFCorpus MedCPT-CE row, which lifts the strongest
biomedical pipeline to nDCG@10 = 0.3664.

The ArguAna collapse of $-0.2428$ is the largest single negative
$\Delta$ in the report and deserves a closer look. A plausible
hypothesis is task / domain mismatch: ArguAna asks for
counter-arguments, where ``relevant'' means \emph{contradicts} the
query, and an entailment-trained cross-encoder may produce misleading
scores on this objective. The competing hypothesis is that the
collapse is a candidate-depth artefact: at $k=100$ the reranker has
many wrong candidates to bring forward, while at smaller $k$ it would
have less room to do damage. Section~\ref{subsec:reranker_depth_ci}
adds a depth ablation and paired bootstrap CIs that let us tell these
two hypotheses apart.

\subsection{Reranker depth and paired confidence intervals}
\label{subsec:reranker_depth_ci}

The $\Delta$ values in Table~\ref{tab:reranker} are point estimates at
candidate depth $k = 100$. Two pieces of evidence are added here to
support (or retract) the reranker claims in
Section~\ref{subsec:reranker}: a candidate-depth ablation at
$k \in \{10, 20, 50, 100\}$, and a paired bootstrap 95\% CI on
$\Delta = \text{nDCG@10}_{\text{rerank}} - \text{nDCG@10}_{\text{first-stage}}$
at the same $k=100$ depth used in Table~\ref{tab:reranker}.

\paragraph{Depth ablation.}
For each (dataset, first-stage) pair the BGE reranker is applied to
the top-$k$ candidates of the first stage, for $k$ in
$\{10, 20, 50, 100\}$. nDCG@10 is then measured on the reranked
list. Table~\ref{tab:reranker_depth} reports $\Delta$ at each depth.
The diagnostic question is whether the negative $\Delta$ values at
$k=100$ in Section~\ref{subsec:reranker} persist at smaller $k$. If
they do, the reranker is fighting the task and a task / domain
mismatch claim is justified. If $\Delta$ moves towards zero (or
becomes positive) at $k=10$, the negative result is a recall-ceiling
artefact rather than a real mismatch, and the claim should be
retracted.

\begin{table}[ht]
    \centering
    \caption{Candidate-depth ablation for the BGE-reranker-base over
    the six (dataset, first-stage) pairs of Section~\ref{subsec:reranker}.
    Cells show $\Delta$ nDCG@10. Empty cells require the Colab rerun
    in \texttt{phase\_b.rerank\_depth\_and\_ci}; the $k=100$ column is
    the same value as in Table~\ref{tab:reranker}.}
    \label{tab:reranker_depth}
    \small
    \begin{tabular}{llSSSS}
\toprule
\textbf{Dataset} & \textbf{First-stage} & \textbf{$k=10$} & \textbf{$k=20$} & \textbf{$k=50$} & \textbf{$k=100$} \\
\midrule
NFCorpus & BGE-small & {---} & {---} & {---} & -0.0386 \\
NFCorpus & BM25 & {---} & {---} & {---} & +0.0006 \\
BioASQ-subset$^{\dagger}$ & BM25 & {---} & {---} & {---} & -0.0464 \\
TREC-COVID & E5-small & {---} & {---} & {---} & +0.0302 \\
SciFact & BGE-small & {---} & {---} & {---} & -0.0109 \\
ArguAna & BGE-small & {---} & {---} & {---} & -0.2428 \\
\bottomrule
\end{tabular}

% Cells show $\Delta$ nDCG@10 = reranked $-$ first-stage. Empty cells await the Colab rerun.
% $^{\dagger}$ degenerate corpus, see BioASQ caveat.

\end{table}

\paragraph{Paired confidence intervals.}
Table~\ref{tab:reranker_ci} reports a paired bootstrap CI on $\Delta$
at $k=100$ for the same six pairs. The bootstrap resamples the test
queries $B = 10{,}000$ times, evaluates both first-stage and reranked
rankings on the same resampled queries, and reports the 2.5\% and
97.5\% quantiles of the resampled mean difference. A CI that strictly
excludes zero is what would let us claim that the reranker
significantly helps or hurts on a given dataset at the 5\% level. A
CI that includes zero means the observed $\Delta$ is consistent with
sampling noise on that test set.

\begin{table}[ht]
    \centering
    \caption{Paired bootstrap 95\% CI for the BGE-reranker $\Delta$ at
    $k=100$. $\hat{\Delta}$ matches the corresponding row in
    Table~\ref{tab:reranker}. The CI and $P(\text{rerank}>\text{FS})$
    columns require the Colab rerun.}
    \label{tab:reranker_ci}
    \small
    \begin{tabular}{llSSSc}
\toprule
\textbf{Dataset} & \textbf{First-stage} & \textbf{$n$} & \textbf{$\hat{\Delta}$} & \textbf{95\% CI} & \textbf{$P(\text{rerank}{>}\text{FS})$} \\
\midrule
NFCorpus & BGE-small & 323 & -0.0386 & {---} & {---} \\
NFCorpus & BM25 & 323 & +0.0006 & {---} & {---} \\
BioASQ-subset$^{\dagger}$ & BM25 & 500 & -0.0464 & {---} & {---} \\
TREC-COVID & E5-small & 50 & +0.0302 & {---} & {---} \\
SciFact & BGE-small & 300 & -0.0109 & {---} & {---} \\
ArguAna & BGE-small & 1,406 & -0.2428 & {---} & {---} \\
\bottomrule
\end{tabular}

% Paired bootstrap, $B = 10{,}000$, at the original $k=100$ candidate depth. CI and $P$ columns await the Colab rerun.
% $^{\dagger}$ degenerate corpus, see BioASQ caveat.

\end{table}

\paragraph{Reading the result.}
The point estimates of $\Delta$ at $k=100$ are already in
Tables~\ref{tab:reranker_depth} and~\ref{tab:reranker_ci}. The
small-$k$ columns and the CI columns are populated by the function
\texttt{phase\_b.rerank\_depth\_and\_ci} on a Colab T4 (the
\texttt{tables/} CSVs are then refreshed via \texttt{python
scripts/build\_tables.py --refresh}). With those numbers in hand, the
ArguAna task-mismatch claim of Section~\ref{subsec:reranker} stays as
written only if the ArguAna row stays negative across all $k$. If
$\Delta$ trends towards zero at $k=10$, the claim is downgraded to a
recall-ceiling observation in the conclusion. Reporting the protocol
and the CSV stubs now makes the next iteration of this report a
one-command refresh rather than a re-edit.

\subsection{Paired bootstrap confidence intervals}
\label{subsec:paired_bootstrap}

The main inferential quantity is the paired difference in per-query
nDCG@10 between two retrievers:
\begin{equation}
    \hat{\Delta}_{A-B} \;=\; \frac{1}{n}\sum_{q \in \mathcal{Q}_{\text{test}}}
    \bigl[\text{nDCG@10}_A(q) - \text{nDCG@10}_B(q)\bigr].
    \label{eq:paired_delta}
\end{equation}
The 95\% CI is estimated with a paired bootstrap of $B = 10{,}000$
resamples of the 323-query NFCorpus test set; both retrievers are
evaluated on the same resampled queries.

\begin{table}[ht]
    \centering
    \caption{Paired bootstrap 95\% CI for pairwise nDCG@10 differences on
    NFCorpus. The first two rows use a paired bootstrap on the fixed
    323-query test set ($B=10{,}000$). The next two use query-subset
    resampling (2{,}000 trials of 323 queries drawn from the
    3{,}237-query union of train, dev and test). $P(A>B)$ is the share
    of resamples in which $A$ has higher mean nDCG@10 than $B$.}
    \label{tab:paired_boot}
    \small
    \begin{tabular}{llSSc}
\toprule
\textbf{Contrast} & \textbf{Scheme} & \textbf{Mean $\hat{\Delta}$} & \textbf{95\% CI} & $P(A{>}B)$ \\
\midrule
Dense $-$ BM25 & Paired bootstrap (test) & +0.0221 & {[-0.0016,\, +0.0449]} & n/a \\
Hybrid $-$ Dense & Point estimate (test) & +0.0208 & {---} & --- \\
Dense $-$ BM25 & Query-subset (3237) & +0.0261 & {[+0.0041,\, +0.0471]} & 99.0\% \\
Hybrid $-$ Dense & Query-subset (3237) & +0.0230 & {[+0.0123,\, +0.0344]} & 100.0\% \\
\bottomrule
\end{tabular}

\end{table}

For Dense vs.\ BM25 the paired bootstrap on the test set gives
$\hat{\Delta}=+0.0221$ with CI $[-0.0016,+0.0449]$. The interval barely
includes zero, so the Dense advantage on this test set is not
conclusive at the 5\% level (it is significant near 7\%). The
query-subset bootstrap from a 3{,}237-query pool gives $+0.0261$ with
CI $[+0.0041,+0.0471]$, strictly above zero. Both schemes agree on the
sign; the second is tighter because it draws from a larger pool. The
difference in width shows how the size of the test set limits effect
estimation at small $n$.

\paragraph{Why a paired bootstrap is preferred over Wilcoxon.}
The Wilcoxon signed-rank test for Dense vs.\ BM25 returns $W=9962.5$,
$p=0.0712$. Wilcoxon ignores magnitude, which is a problem in
retrieval: about 33\% of NFCorpus queries are total failures (loss
$=1$) for both systems and contribute zero signed rank. Head-to-head
counts are: BM25 wins 97 queries, Dense wins 118, and 108 are ties.
The paired bootstrap uses the magnitudes and is the appropriate tool
for an effect-size claim such as ``Dense beats BM25 by about 2 nDCG@10
points''. The disagreement between the two tests is informative:
Dense's advantage comes from the size of the wins, not from the number
of wins. The per-query loss correlation between BM25 and Dense is
$r = 0.768$, which is the regime in which paired resampling has more
power than Wilcoxon.

\subsection{Regression for the per-query advantage}
\label{subsec:regression}

The previous round of the project used stratified bins on query
features to study conditional heterogeneity in Dense's advantage.
Binning requires choosing thresholds and discards within-bin variance.
A direct regression replaces the binning:
\begin{equation}
    \Delta(q) \;=\; \beta_0
    + \beta_{\text{gap}} \cdot \widetilde{\text{gap}}(q)
    + \beta_{\text{len}} \cdot \widetilde{\text{len}}(q)
    + \beta_{\text{tech}} \cdot \widetilde{\text{tech}}(q)
    + \varepsilon,
    \label{eq:delta_regression}
\end{equation}
where $\Delta(q) = \text{nDCG@10}_{\text{Dense}}(q) -
\text{nDCG@10}_{\text{BM25}}(q)$ and the tilde indicates $z$-score
normalisation. The vocabulary gap is
$1 - \max_{d \in \text{qrel}(q)} \text{overlap}(q, d)$; query length is
the number of content tokens; technicality is the mean IDF of content
tokens. Coefficients are estimated by OLS on $n=323$ test queries
(Table~\ref{tab:regression}). 95\% CIs are bootstrap percentile
intervals with $B=10{,}000$.

\begin{table}[ht]
    \centering
    \caption{OLS regression of $\Delta(q) = $ Dense $-$ BM25 per-query
    nDCG@10 on standardised query features. All features are $z$-scored,
    so coefficient sizes are comparable.}
    \label{tab:regression}
    \small
    \begin{tabular}{lSSSc}
\toprule
\textbf{Feature} & \textbf{Coef} & \textbf{95\% CI} & \textbf{Boot $p$} & \textbf{Sig.} \\
\midrule
Intercept & +0.0233 & {[+0.0014,\, +0.0459]} & 0.037 & $*$ \\
Vocabulary gap & +0.0475 & {[+0.0249,\, +0.0713]} & {$<\!0.001$} & $***$ \\
Query length (tokens) & +0.0112 & {[-0.0168,\, +0.0403]} & 0.430 & {---} \\
Technicality (mean IDF) & -0.0343 & {[-0.0648,\, -0.0053]} & 0.019 & $*$ \\
\midrule
$R^2$ & 0.071 & & & \\
$n$ & 323 & & & \\
\bottomrule
\end{tabular}

\end{table}

Two coefficients are significant. A one-standard-deviation increase in
vocabulary gap raises Dense's advantage by $+0.0475$ nDCG points; this
is the regime in which BM25 cannot match because the relevant document
shares little surface vocabulary with the query. A one-SD increase in
mean query IDF reduces Dense's advantage by $-0.0343$ nDCG points, so
on technical queries the general-domain dense encoder loses ground.
Query length is not significant after the other two features are
controlled for. Bootstrap Spearman correlations confirm the regression
direction: vocabulary gap $\rho=+0.199$ ($p=3\!\times\!10^{-4}$), query
length $\rho=+0.230$, and technicality $\rho=-0.150$. The $R^2$ of
0.071 is small. Three query features explain about 7\% of $\Delta(q)$
variance; most of the rest reflects per-query factors (including
document-side properties such as how well the corpus covers a topic)
that would require a separate model.

\subsection{Per-query router}
\label{subsec:router}

The regression in Equation~\ref{eq:delta_regression} can be the first
stage of an algorithmic decision rule. A per-query router predicts
which retrieval strategy to use, choosing among four candidates:
\begin{itemize}[noitemsep,topsep=2pt]
    \item BM25 (lexical sparse, no compute on the query).
    \item BGE-small (the strongest single dense bi-encoder on NFCorpus).
    \item Static linear hybrid of BM25 and BGE-small at $\alpha=0.40$.
    The router treats this as one discrete option; a richer per-query
    $\alpha(q)$ regression that varies the weight continuously is left
    to future work, because it would lose the well-defined routing
    label space $\arg\max_s \text{nDCG@10}_s(q)$ and would make the
    oracle upper bound infinite.
    \item BGE-small as first stage with a BGE cross-encoder reranker
    over the top-100 candidates.
\end{itemize}

\paragraph{Features.}
Six per-query features are computed: vocabulary gap $g_q$, content-token
length $\ell_q$, mean IDF $\tau_q$, a binary question-form flag (regex
match against question words), the share of medical-vocabulary tokens
(a UMLS-lite affix list), and the BM25 top-1 score (a cheap easy-query
signal). The first three are the same features as in
Section~\ref{subsec:regression}; the last three are added because
$R^2 = 0.071$ leaves substantial $\Delta(q)$ variance unexplained.

\paragraph{Disclaimer: one feature is retrospective.}
The vocabulary gap is computed against the ground-truth relevant
documents (the qrels). At deployment time the qrels are not known, so
this feature is not available to a real router. The router results in
this section therefore include a retrospective signal and are best
read as a near-oracle upper bound on what a router with this feature
family could reach, not as a deployment number. The other five
features (query length, mean IDF, question-form flag, medical-token
share, BM25 top-1 score) are computable from the query alone or from
cheap first-stage scores, so they are deployable. We keep the
retrospective feature on purpose. The main finding below is negative
(the learned routers do not beat the static hybrid), and that finding
becomes stronger rather than weaker when the routers fail with the
retrospective feature on board. Replacing the vocabulary-gap feature
with deployable proxies (e.g.\ BM25/Dense rank agreement, score margins)
is left to future work.

\paragraph{Models.}
Three router variants are trained on the NFCorpus train and dev splits
with per-query labels $\arg\max_s \text{nDCG@10}_s(q)$ over the four
strategies $s$: a multinomial logistic regression on $z$-scored
features, a LightGBM multi-class classifier with 300 trees, and an
oracle that picks the per-query best of four. The oracle is not
deployable; it is reported as the upper bound any router with this
strategy set can reach.

\begin{table}[ht]
    \centering
    \caption{Query-routing evaluation on the NFCorpus test split. The
    Always-baselines reproduce the strongest single retriever and the
    static hybrid; the Always BGE-small + BGE-reranker baseline comes
    from Section~\ref{subsec:reranker}; the Oracle row is the per-query
    best-of-four upper bound.}
    \label{tab:router}
    \small
    \begin{tabular}{lSS}
\toprule
\textbf{Method} & \textbf{nDCG@10} & \textbf{$\Delta$ vs best Always} \\
\midrule
Always BM25 & 0.3076 & -0.0358 \\
Always BGE-small & 0.3393 & -0.0041 \\
Static Hybrid (BM25 $+$ BGE-small, $\alpha\!=\!0.40$) & \textbf{0.3434} & 0.0000 \\
Always BGE-small $+$ BGE-reranker@100 & 0.3007 & -0.0427 \\
\midrule
Logistic router & 0.3230 & -0.0204 \\
LightGBM router & 0.3325 & -0.0109 \\
\midrule
Oracle router (upper bound) & 0.4097 & +0.0663 \\
\bottomrule
\end{tabular}

\end{table}

\paragraph{Hybrid composition.}
The Static Hybrid in Table~\ref{tab:router} fuses BM25 with
\emph{BGE-small} at $\alpha = 0.40$, and is the same hybrid the router
sees as one of its four candidates. This is different from the hybrid
of Section~\ref{subsec:complementarity}, which fuses BM25 with
\emph{MiniLM-Dense} at $\alpha = 0.40$ tuned on the NFCorpus dev split
(value 0.3380). The router uses BGE-small inside the hybrid because
otherwise Hybrid would be a duplicate of Always BGE-small. The
0.0054 difference between the two hybrid numbers is a definition
difference in the dense component, not a re-tuning of $\alpha$.

\paragraph{Reading the result.}
The main finding is negative. Even with the retrospective
vocabulary-gap feature on board, the learned routers do not beat the
strongest static baseline. LightGBM reaches 0.3325 and the logistic
router 0.3230 on the test split; the Static Hybrid reaches 0.3434
and Always BGE-small reaches 0.3393. Both learned routers do beat
Always BM25 ($+0.025$ for LightGBM, $+0.015$ for logistic), so the
routing signal is real, but it is smaller than the gain from picking
the right static fusion in the first place.

The Oracle row paints a different picture. The per-query best-of-four
reaches 0.4097 on the test split, which is $+0.0663$ above Static
Hybrid. The per-query heterogeneity is therefore large: a perfect
router with this strategy set would lift the test set by 19\% relative
to the best static baseline. The interesting question is not whether
routing is worth doing in principle but why a simple learned router
fails to close most of this oracle gap. Two reasons are visible. The
six features explain only about 7\% of $\Delta(q)$ variance
(Section~\ref{subsec:regression}, $R^2 = 0.071$); most of the
per-query variance lives in signals these features do not capture.
The label space is also small: the router chooses among four discrete
strategies, which limits how much fine-grained per-query control any
classifier can deliver. A continuous per-query mixing weight
$\alpha(q)$ would give the oracle more room and would likely give a
learned router more to fit, at the cost of making the oracle upper
bound less well-defined.

The LightGBM feature-importance gain is reported in
Table~\ref{tab:router_features}. The dominant predictor is
\emph{bm25\_top1}, the BM25 top-1 score, a cheap easy-query signal
that overwhelms the more interpretable vocabulary-gap and technicality
features. \emph{tech} is a close second; the question-form flag and
medical-vocabulary share are negligible. The router therefore
exploits query difficulty (proxied by BM25 top-1) and technical
density rather than the lexical-overlap gap that the descriptive
regression highlighted.

\begin{table}[ht]
    \centering
    \caption{LightGBM router feature importance (gain).}
    \label{tab:router_features}
    \small
    \begin{tabular}{lSS}
\toprule
\textbf{Feature} & \textbf{Gain} & \textbf{Share of total} \\
\midrule
bm25\_top1 & 15,387.6 & 38.7\% \\
tech & 14,542.3 & 36.6\% \\
gap & 6,488.0 & 16.3\% \\
len & 2,718.8 & 6.8\% \\
med\_share & 530.3 & 1.3\% \\
is\_question & 358.0 & 0.9\% \\
\bottomrule
\end{tabular}

\end{table}

The router is the only model in the study that is fit to the NFCorpus
query distribution. Its train--test gap therefore provides the only
non-trivial overfitting signal, discussed in Section~\ref{subsec:gap}.

\subsection{Practical efficiency}
\label{subsec:efficiency}

Retrieval systems are also judged by their efficiency.
Table~\ref{tab:efficiency} reports indexing time, per-query latency
(single-threaded CPU), throughput and index memory on the
3{,}633-document NFCorpus corpus.

\begin{table}[ht]
    \centering
    \caption{Efficiency metrics on NFCorpus (3{,}633 documents,
    single-threaded, top-100 retrieval). Index build includes corpus
    encoding for dense methods. Throughput is measured on the 323-query
    test set. Memory refers to the in-memory index size: inverted list
    for BM25, scipy.sparse CSR for SPLADE, and float32 matrix for dense
    encoders. Dense encoder build times use a Colab T4 GPU; CPU-only
    builds would be roughly $10\times$ slower.}
    \label{tab:efficiency}
    \small
    \begin{tabular}{lrrrl}
\toprule
\textbf{Retriever} & \textbf{Index build (s)} & \textbf{Latency (ms/q)} & \textbf{Throughput (q/s)} & \textbf{Index memory} \\
\midrule
BM25 & 12.3 & 2.4 & 412 & inverted list (sparse) \\
Dense (MiniLM) & 11.6 & 8.6 & 116 & 5 MB (3633 x 384 f32) \\
BGE-small & 32.2 & 0.7 & 1,389 & 5 MB (3633 x 384 f32) \\
E5-small & 33.1 & 0.6 & 1,561 & 5 MB (3633 x 384 f32) \\
SPLADE & 38.4 & 2.8 & 353 & sparse CSR, ~151 nnz/doc \\
MedCPT & 51.4 & 1.2 & 867 & 11 MB (3633 x 768 f32) \\
\bottomrule
\end{tabular}

\end{table}

BGE-small is at the same time the strongest bi-encoder by quality and
the fastest at query time (1{,}389 q/s). SPLADE matches dense quality
while indexing into a sparse matrix that scales better with corpus
size. MedCPT pays the cost of a dual-encoder and a 768-dimensional
representation: twice the memory of the 384-dimensional encoders and a
second BERT-base forward pass per query. BM25 is the cheapest to
index (12.3\,s, no neural forward pass). On a corpus of this size
every retriever runs in real time on CPU; the trade-offs become
relevant on larger corpora and memory-constrained deployments.

\subsection{How many queries are needed to distinguish systems?}
\label{subsec:n_required}

For a bounded loss in $[0,1]$, a Hoeffding-style requirement gives the
sample size needed to certify a pairwise gap $\hat{\Delta}_{A-B}$ at
confidence $1-\delta = 0.95$:
\[
n_{\min} \ge \left\lceil \frac{\log(2/\delta)}{2\hat{\Delta}_{A-B}^2} \right\rceil.
\]
Table~\ref{tab:n_required} reports $n_{\min}$ for four pairwise
contrasts of interest.

\begin{table}[ht]
    \centering
    \caption{Pairwise sample-size requirements (Hoeffding-style,
    $\delta=0.05$) using observed NFCorpus nDCG@10 gaps.}
    \label{tab:n_required}
    \small
    \begin{tabular}{lSS}
\toprule
\textbf{Pair} & \textbf{Observed gap $\hat{\Delta}$} & \textbf{$n_{\min}$} \\
\midrule
BGE-small $-$ BM25 & +0.0448 & 919 \\
BGE-small $-$ E5-small & +0.0121 & 12,598 \\
Hybrid $-$ Dense & +0.0208 & 4,274 \\
Dense $-$ BM25 & +0.0233 & 3,394 \\
\bottomrule
\end{tabular}

\end{table}

Each row is tied to a specific claim. Row 1 corresponds to the main
NFCorpus ranking (Tables~\ref{tab:nfcorpus_6way}
and~\ref{tab:multi_dataset}); the 919-query requirement says the
existing 323-query test set is too small to certify the gap at 95\%
under the conservative bound, but the paired bootstrap of
Section~\ref{subsec:paired_bootstrap} already covers it with much
higher precision. Row 2 (BGE-small $-$ E5-small) is the within-bi-encoder
gap of Section~\ref{subsec:nf_quality}; distinguishing the two
strongest dense encoders at 95\% requires an order of magnitude more
queries. Row 3 (Hybrid $-$ Dense) is the complementarity gap of
Section~\ref{subsec:complementarity}, bounded empirically by the
query-subset CI $[+0.0123,\,+0.0344]$ in Table~\ref{tab:paired_boot}.
Row 4 (Dense $-$ BM25) is the contrast that drives the regression and
the router; significant under query-subset resampling but not
certifiable by Hoeffding on $n=323$.

The same formula can be applied to the reranker effects of
Table~\ref{tab:reranker}: for BGE-reranker on TREC-COVID
($\hat\Delta = +0.0302$), $n_{\min} = 2{,}023$; for MedCPT-CE on
BioASQ-subset ($+0.0513$), $n_{\min} = 701$; for MedCPT-CE on NFCorpus
($+0.0271$), $n_{\min} = 2{,}512$; for the BGE-reranker collapse on
ArguAna ($-0.2428$), $n_{\min} = 32$. The asymmetry between certifying
small in-domain gains (uncertifiable on the available test sets) and
certifying large out-of-domain failures (trivially certifiable) is
informative on its own.

\subsection{Query-subsampling sensitivity}

To check whether the ranking is stable under different query samples,
the nDCG@10 estimate is repeated over 2{,}000 random subsamples of 323
queries each, drawn with replacement from the union of train, dev and
test (3{,}237 queries in total).

\begin{table}[ht]
    \centering
    \caption{Subsampling sensitivity across 2{,}000 bootstrap trials.
    Dense $-$ BM25 has a lower 2.5th percentile strictly above zero,
    so Dense wins in at least 97.5\% of resamples.}
    \label{tab:subsampling}
    \small
    \begin{tabular}{lSSSS}
\toprule
\textbf{Quantity} & \textbf{Mean} & \textbf{$Q_{2.5}$} & \textbf{$Q_{97.5}$} & \textbf{Std} \\
\midrule
BM25 & 0.2924 & 0.2610 & 0.3238 & 0.0160 \\
Dense & 0.3185 & 0.2865 & 0.3514 & 0.0164 \\
Hybrid & 0.3415 & 0.3089 & 0.3741 & 0.0166 \\
\midrule
Dense $-$ BM25 & 0.0261 & 0.0041 & 0.0471 & 0.0111 \\
Hybrid $-$ Dense & 0.0230 & 0.0123 & 0.0344 & 0.0056 \\
\bottomrule
\end{tabular}

\end{table}

\subsection{Train--test generalisation gap}
\label{subsec:gap}

The router of Section~\ref{subsec:router} is the only model in the
study that is fit to NFCorpus queries, so the train--test gap
discussion is anchored on it. The four first-stage retrievers (BM25,
Dense, BGE-small, E5-small) are not fit; their train--test difference
is sampling noise. Table~\ref{tab:gap} gives
$\Delta_{\text{gen}} = R_{\text{test}} - R_{\text{train}}$.

\begin{table}[ht]
    \centering
    \caption{Train--test generalisation gap on NFCorpus. The four
    first-stage retrievers are not fit to NFCorpus queries, so their
    gap is the variance baseline; the routers are fit. The
    Mann--Whitney $p$-value tests whether per-query loss distributions
    differ between the train split (2{,}590 queries) and the test split
    (323 queries).}
    \label{tab:gap}
    \small
    \begin{tabular}{lSSSSSS}
\toprule
\textbf{Quantity} & \textbf{BM25} & \textbf{Dense} & \textbf{BGE-small} & \textbf{E5-small} & \textbf{Logistic} & \textbf{LightGBM} \\
\midrule
$R_{\text{train}}$ & 0.7047 & 0.6792 & 0.6615 & 0.6743 & 0.6841 & 0.6207 \\
$R_{\text{test}}$ & 0.7060 & 0.6827 & 0.6612 & 0.6733 & 0.6770 & 0.6675 \\
Gap $\Delta_{\text{gen}}$ & +0.0013 & +0.0035 & -0.0003 & -0.0009 & -0.0071 & \textbf{+0.0468} \\
Mann--Whitney $p$ & 0.9252 & 0.8768 & 0.9109 & 0.8867 & {---} & {---} \\
\bottomrule
\end{tabular}

\end{table}

The four first-stage gaps are all below $0.004$ in absolute value and
all four Mann--Whitney $p$-values exceed $0.87$: the train and test
loss distributions are statistically indistinguishable, and any gap
below this threshold is sampling noise on this dataset. SPLADE and
MedCPT are not included because encoding the 2{,}590 training queries
through their full pipelines would have required substantially more
GPU time.

The two router columns provide the actual generalisation comparison.
The logistic router has $R_{\text{train}} = 0.6841$ and
$R_{\text{test}} = 0.6770$, gap $-0.0071$; the negative sign is
sampling noise and the magnitude is similar to the first-stage
baseline. The linear model with six $z$-scored features is too
constrained to overfit 2{,}914 training queries. The LightGBM router,
with 300 trees, has $R_{\text{train}} = 0.6207$ and $R_{\text{test}} =
0.6675$, gap $+0.0468$; this is more than ten times the
variance-baseline threshold. With 300 trees on six features and
2{,}914 queries, the boosted model fits training-query idiosyncrasies
that do not generalise. This is the regime that the
train--test gap framework is designed to detect: a non-fitted
estimator should stay at the variance baseline, a fitted estimator
should track its capacity. Both predictions hold. The LightGBM router
also illustrates the cost of the extra capacity: it beats the logistic
router by $+0.0095$ on the test set but does not beat the static
hybrid baseline. A capacity-controlled router (e.g., LightGBM with 30
trees, or the linear model with cross-validated $L_2$) is the natural
next step.

% -------------------------------------------------------
\section{Discussion}
\label{sec:theory}

\subsection{No Free Lunch and inductive-bias alignment}

The No Free Lunch theorem says that no learning algorithm dominates
across all distributions. In retrieval the equivalent statement is
that performance depends on the alignment between the retriever's
inductive bias and the structure of query--document pairs. The
multi-dataset sweep of Section~\ref{subsec:multi} tests this:
no single retriever is best on every dataset. BGE-small wins three
benchmarks and is second on a fourth; E5-small wins TREC-COVID; BM25
wins the no-distractor BioASQ-subset (with the caveat above). The
$\Delta(q)$ regression of Section~\ref{subsec:regression} makes the
same point inside one dataset: vocabulary gap and technicality are
signed predictors of which queries favour Dense over BM25.

\subsection{Empirical vs.\ true risk}
\label{subsec:risk}

No first-stage retriever in this study is fitted to NFCorpus queries,
so the train--test split reflects evaluation variance rather than
overfitting. The small first-stage gaps in Table~\ref{tab:gap}
($|\Delta_{\text{gen}}| \le 0.004$) and the paired-bootstrap CIs in
Table~\ref{tab:paired_boot} together describe how well $\hat{R}$
approximates $R$. With $n=323$, the 95\% bootstrap CI half-width for
Dense is about 0.034: the true risk is known to within roughly $\pm 3$
nDCG@10 points. The contrast between this width and the
distribution-free Hoeffding width (0.076) is the practical gain from a
variance-aware estimator, a roughly $2.3\times$ improvement.

\subsection{Stability}

Two stability questions are tested. \emph{Sample stability}: under
2{,}000 query-subset resamples, Dense beats BM25 in about 99\% of
resamples, and Hybrid beats Dense in about 100\%; on NFCorpus the
ranking BM25 $<$ Dense $<$ Hybrid is highly stable.
\emph{Cross-dataset stability}: the same ranking does not hold across
datasets. On TREC-COVID, BM25 beats MiniLM-Dense by 0.122 nDCG@10. On
ArguAna, MedCPT drops below BM25 by 0.228. Sample stability inside a
dataset is high; cross-dataset stability is low. This is the expected
signature of an inductive bias whose alignment with the data is
distribution-dependent.

\subsection{Complementarity and hybrid fusion}
\label{subsec:complementarity}

Tuning $\alpha$ on the NFCorpus dev split for the BM25--MiniLM hybrid
of Equation~\ref{eq:hybrid} gives $\alpha^\star = 0.40$. Applied on
the test split the hybrid reaches nDCG@10 $= 0.3380$, which is
$+0.0208$ above MiniLM-Dense alone and $+0.0441$ above BM25. RRF with
$k=60$ reaches 0.3235, which improves over each base system but is
weaker than tuned linear fusion. The query-subset CI for Hybrid $-$
Dense is $[+0.0123, +0.0344]$ with $P(\text{Hybrid}>\text{Dense})=100\%$,
so fusion consistently helps.

The Pearson correlation of per-query BM25 and Dense losses is
$r = 0.768$: high but not perfect. This is the regime that score
fusion is useful in. Stratifying by vocabulary gap and by technicality
(Tables~\ref{tab:vocab_gap_strat} and~\ref{tab:tech_strat}) clarifies
when fusion helps. On low-gap queries, where the relevant document
already shares most vocabulary with the query, BM25 beats Dense by
$-0.034$ and fusion adds $+0.065$ over Dense. On high-gap queries
Dense beats BM25 by $+0.067$, and fusion loses slightly because the
exact-match signal is not useful. Technicality behaves symmetrically:
on highly technical queries, Dense loses to BM25 by $-0.020$ and
fusion adds $+0.049$ over Dense by re-introducing exact matching.

\begin{table}[ht]
\centering
\caption{Vocabulary-gap stratification of NFCorpus test (tertile split).
The vocabulary gap is $1 - \max_{d \in \text{qrel}(q)}
\text{overlap}(q, d)$.}
\label{tab:vocab_gap_strat}
\small
\begin{tabular}{lSSSSS}
\toprule
\textbf{Stratum} & \textbf{$n$} & \textbf{BM25} & \textbf{Dense} & \textbf{Hybrid} & \textbf{Dense$-$BM25} \\
\midrule
Low gap & 108 & 0.5267 & 0.4931 & \textbf{0.5577} & -0.0336 \\
Medium gap & 107 & 0.3039 & 0.3406 & \textbf{0.3580} & +0.0367 \\
High gap & 108 & 0.0513 & \textbf{0.1183} & 0.0986 & +0.0670 \\
\bottomrule
\end{tabular}

\end{table}

\begin{table}[ht]
\centering
\caption{Technicality (mean IDF) stratification of NFCorpus test
(tertile split).}
\label{tab:tech_strat}
\small
\begin{tabular}{lSSSSS}
\toprule
\textbf{Stratum} & \textbf{$n$} & \textbf{BM25} & \textbf{Dense} & \textbf{Hybrid} & \textbf{Dense$-$BM25} \\
\midrule
Low IDF / plainer & 108 & 0.3193 & 0.3532 & \textbf{0.3570} & +0.0338 \\
Mid IDF & 107 & 0.3124 & 0.3690 & \textbf{0.3789} & +0.0566 \\
High IDF / technical & 108 & 0.2503 & 0.2301 & \textbf{0.2786} & -0.0202 \\
\bottomrule
\end{tabular}

\end{table}

The vocabulary-gap Spearman correlation with $\Delta(q)$ is
$\rho_{\text{gap}} = +0.199$ ($p = 3\!\times\!10^{-4}$): Dense's
per-query advantage grows as the lexical overlap with relevant
documents shrinks. The equivalent overlap correlation is $-0.199$
(since $\text{gap} = 1 - \text{overlap}$); the report uses the gap
form so that the sign matches $\beta_{\text{gap}} > 0$ in the
regression.

% -------------------------------------------------------
\section{Downstream Generation: MIRAGE/MedRAG}
\label{sec:downstream}

To ground the RAG framing of the report, a compact downstream
generation experiment is added on a biomedical QA task with gold
answers~\cite{Xiong2024MIRAGE}. The task is
PubMedQA-labeled~\cite{Jin2019PubMedQA} (\textit{yes/no/maybe}
classification); BioASQ Y/N~\cite{Tsatsaronis2015BioASQ} was planned but
not produced on this run (see table note). The pipeline retrieves
top-5 evidence passages with three retrievers (BM25, BGE-small,
BGE-small + BGE-reranker), prompts the generator with question and
evidence, and scores the prediction against the gold label. A
no-retrieval closed-book baseline is included. Two generators were
planned: \texttt{google/flan-t5-base} and
\texttt{meta-llama/Llama-3.2-3B-Instruct} (4-bit). Only flan-t5-base
was actually run (Llama is gated and the Colab session did not
include the required login).

\begin{table}[ht]
    \centering
    \caption{Downstream MIRAGE/MedRAG accuracy on the PubMedQA-labeled
    test set ($n = 1{,}000$). The closed-book row ablates the retrieval
    contribution. Llama-3.2-3B-Instruct rows and the BioASQ Y/N task
    were planned but are not reported in this draft (see note below).}
    \label{tab:mirage}
    \small
    \begin{tabular}{lllS}
\toprule
\textbf{Task} & \textbf{Retriever} & \textbf{Generator} & \textbf{Accuracy} \\
\midrule
PubMedQA-labeled & None (closed-book) & flan-t5-base & 0.379 \\
PubMedQA-labeled & BM25 & flan-t5-base & \textbf{0.561} \\
PubMedQA-labeled & BGE-small & flan-t5-base & \textbf{0.563} \\
PubMedQA-labeled & BGE-small + BGE-reranker & flan-t5-base & \textbf{0.561} \\
\midrule
\multicolumn{4}{l}{\emph{Llama-3.2-3B-Instruct rows: not run (gated; no HF login in this Colab session).}} \\
\multicolumn{4}{l}{\emph{BioASQ Y/N rows: not run (loader returned no yes/no records).}} \\
\bottomrule
\end{tabular}

\end{table}

\paragraph{Read this section as a saturation sanity check.}
The PubMedQA-labeled experiment is included to confirm that the
pipeline carries evidence end-to-end, not as evidence about
retrieval-method choice for RAG. The PubMedQA-labeled corpus ships
only ten candidate passages per question, and the gold passage is
almost always inside that pool. Once the top-five evidence reaches
the generator, all three retrievers (BM25, BGE-small, BGE-small +
BGE-reranker) feed essentially the same passages. The point estimates
in Table~\ref{tab:mirage} are 0.561, 0.563 and 0.561 -- within
$\pm 0.002$ and well below $1\sigma$ of a binomial at $p=0.56$
(about 0.016). The retrieval comparison is therefore uninformative
by construction; it is reported only for completeness and is
\emph{not} used in this report as evidence that BM25, BGE-small or
the BGE reranker are equivalent on RAG tasks in general.

One observation is worth keeping. Closed-book flan-t5-base reaches
0.379 on this 3-way task (chance is $\approx 0.34$), and adding any
top-5 retrieval lifts it to $\approx 0.56$. The $+0.18$ absolute lift
is not a clean test of retriever quality -- the candidate pool is
too small for that -- but it does confirm the pipeline propagates
retrieved evidence to the generator. A clean RAG retrieval-choice
experiment would rebuild this task with a realistic distractor
corpus (e.g.\ a PubMed-abstracts index at the 100k-document scale).
This is listed in Section~\ref{sec:limitations} and left to future
work.

% -------------------------------------------------------
\section{Limitations}
\label{sec:limitations}

\textbf{(1) Downstream scope is narrow.} Only PubMedQA-labeled with
flan-t5-base is reported; Llama-3.2-3B-Instruct and BioASQ Y/N rows
were planned but not produced.

\textbf{(2) Synonym perturbation excluded.} The midterm report
included a synonym perturbation experiment to measure per-query
stability under small textual changes. It was removed because of a
reproducibility issue on macOS related to the \texttt{faiss-cpu}
library and OpenMP interaction with PyTorch. The subsampling analysis
in Table~\ref{tab:subsampling} partially addresses stability by
varying the query set rather than the query text.

\textbf{(3) CPU-only timing.} All latency and throughput numbers are
single-threaded CPU measurements. Dense encoder throughput would be
much higher with GPU-accelerated FAISS. The conclusions in
Section~\ref{subsec:efficiency} apply to CPU inference, not to
production GPU pipelines.

\textbf{(4) PubMedQA candidate pool is too small to test retrieval.}
The PubMedQA-labeled task in Section~\ref{sec:downstream} ships only
ten candidate passages per question, so all three retrievers feed
near-identical evidence to the generator and the rows in
Table~\ref{tab:mirage} saturate at $\approx 0.56$. The
BM25/BGE-small/+reranker comparison is therefore not interpretable
as RAG retrieval-choice evidence. Rebuilding the task with a
realistic distractor corpus (e.g.\ a 100k-document PubMed-abstracts
index) is the natural fix.

\textbf{(5) BioASQ subset is degenerate on this run.} BioASQ-subset
contains 28{,}001 documents in a qrel-union pool with zero distractors
because the loader that succeeded provided a small public mirror.
This biases the BioASQ row towards lexical retrieval. The row is
reported as a stress / control case only and is not used as
biomedical retrieval evidence (see
Section~\ref{subsec:multi}). Re-running with the gated official
BioASQ-BEIR release (14M documents) or with a custom
PubMed-distractor pool is the natural fix.

% -------------------------------------------------------
\section{Conclusion}

The empirical picture is non-uniform on every axis. No single
retriever wins on all five datasets: BGE-small wins NFCorpus, SciFact
and ArguAna, and is second on TREC-COVID, where E5-small is best.
BioASQ-subset is degenerate on this run and is excluded from
biomedical claims (Section~\ref{subsec:multi}). On the four
non-degenerate datasets BGE-small averages $+0.060$ above BM25, which
is the largest mean of any learned system. MedCPT collapses to 0.1332
on ArguAna ($-0.227$ vs.\ BM25), the largest undisputed negative gap
in the study.

A stronger general-domain reranker does not give a stronger pipeline.
BGE-reranker-base hurts on four of five datasets and only TREC-COVID
gains ($+0.030$). The clean positive case is the in-domain MedCPT
cross-encoder on NFCorpus ($+0.027$), which has the opposite sign of
the first-stage MedCPT result on the same data. Domain alignment is
most useful when the model attends jointly to query and document, not
when query and document are encoded separately.

The per-query router finds real heterogeneity but does not beat the
best static fusion. LightGBM beats Always-BM25 by $+0.025$ but loses
to the static BM25+BGE-small hybrid by $-0.011$, even with a
retrospective vocabulary-gap feature on board (see disclaimer in
Section~\ref{subsec:router}). The oracle is $+0.066$ above the static
hybrid, so the routing signal exists, but the six features explain
only about 7\% of $\Delta(q)$. The train--test gap follows the
expected capacity-vs-sample-size pattern: the linear router stays
near the variance baseline, the boosted router has a gap of $+0.0468$
and still does not beat the static hybrid. A capacity-controlled
router with a fully deployable feature set is the natural next
step.

The aggregate reading is an instance of No Free Lunch. No retriever
dominates, no reranker dominates, and a fitted router does not yet
beat the best static fusion. The most reliable biomedical gain in
this study did not come from any of the four stronger baselines added
to the first-stage suite, but from a domain-aligned cross-encoder on
top of a weaker first stage. The methodological contribution is to
make the heterogeneity story testable: the oracle bound, the
train--test variance baseline, and the Hoeffding $n_{\min}$ table
together turn the conditional-heterogeneity argument into an explicit
sample-size budget.

% -------------------------------------------------------
\printbibliography

\end{document}
